% Options for packages loaded elsewhere
\PassOptionsToPackage{unicode}{hyperref}
\PassOptionsToPackage{hyphens}{url}
\documentclass[
  american,
  11pt,
]{article}
\usepackage{xcolor}
\usepackage[margin=25mm]{geometry}
\usepackage{amsmath,amssymb}
\setcounter{secnumdepth}{-\maxdimen} % remove section numbering
\usepackage{iftex}
\ifPDFTeX
  \usepackage[T1]{fontenc}
  \usepackage[utf8]{inputenc}
  \usepackage{textcomp} % provide euro and other symbols
\else % if luatex or xetex
  \usepackage{unicode-math} % this also loads fontspec
  \defaultfontfeatures{Scale=MatchLowercase}
  \defaultfontfeatures[\rmfamily]{Ligatures=TeX,Scale=1}
\fi
\usepackage{lmodern}
\ifPDFTeX\else
  % xetex/luatex font selection
\fi
% Use upquote if available, for straight quotes in verbatim environments
\IfFileExists{upquote.sty}{\usepackage{upquote}}{}
\IfFileExists{microtype.sty}{% use microtype if available
  \usepackage[]{microtype}
  \UseMicrotypeSet[protrusion]{basicmath} % disable protrusion for tt fonts
}{}
\makeatletter
\@ifundefined{KOMAClassName}{% if non-KOMA class
  \IfFileExists{parskip.sty}{%
    \usepackage{parskip}
  }{% else
    \setlength{\parindent}{0pt}
    \setlength{\parskip}{6pt plus 2pt minus 1pt}}
}{% if KOMA class
  \KOMAoptions{parskip=half}}
\makeatother
% definitions for citeproc citations
\NewDocumentCommand\citeproctext{}{}
\NewDocumentCommand\citeproc{mm}{%
  \begingroup\def\citeproctext{#2}\cite{#1}\endgroup}
\makeatletter
 % allow citations to break across lines
 \let\@cite@ofmt\@firstofone
 % avoid brackets around text for \cite:
 \def\@biblabel#1{}
 \def\@cite#1#2{{#1\if@tempswa , #2\fi}}
\makeatother
\newlength{\cslhangindent}
\setlength{\cslhangindent}{1.5em}
\newlength{\csllabelwidth}
\setlength{\csllabelwidth}{3em}
\newenvironment{CSLReferences}[2] % #1 hanging-indent, #2 entry-spacing
 {\begin{list}{}{%
  \setlength{\itemindent}{0pt}
  \setlength{\leftmargin}{0pt}
  \setlength{\parsep}{0pt}
  % turn on hanging indent if param 1 is 1
  \ifodd #1
   \setlength{\leftmargin}{\cslhangindent}
   \setlength{\itemindent}{-1\cslhangindent}
  \fi
  % set entry spacing
  \setlength{\itemsep}{#2\baselineskip}}}
 {\end{list}}
\usepackage{calc}
\newcommand{\CSLBlock}[1]{\hfill\break\parbox[t]{\linewidth}{\strut\ignorespaces#1\strut}}
\newcommand{\CSLLeftMargin}[1]{\parbox[t]{\csllabelwidth}{\strut#1\strut}}
\newcommand{\CSLRightInline}[1]{\parbox[t]{\linewidth - \csllabelwidth}{\strut#1\strut}}
\newcommand{\CSLIndent}[1]{\hspace{\cslhangindent}#1}
\ifLuaTeX
\usepackage[bidi=basic,shorthands=off]{babel}
\else
\usepackage[bidi=default,shorthands=off]{babel}
\fi
\ifLuaTeX
  \usepackage{selnolig} % disable illegal ligatures
\fi
\setlength{\emergencystretch}{3em} % prevent overfull lines
\providecommand{\tightlist}{%
  \setlength{\itemsep}{0pt}\setlength{\parskip}{0pt}}
\usepackage{bookmark}
\IfFileExists{xurl.sty}{\usepackage{xurl}}{} % add URL line breaks if available
\urlstyle{same}
\hypersetup{
  pdftitle={Local curvature of the nuisance-profile score},
  pdflang={en-US},
  hidelinks,
  pdfcreator={LaTeX via pandoc}}

\title{Local curvature of the nuisance-profile score}
\author{}
\date{}

\begin{document}
\maketitle

\subsection{Incidental initial states: a local bias
calculation}\label{incidental-initial-states-a-local-bias-calculation}

Increasing the number of record-specific initial states raises an
incidental-parameter issue {[}1, 2{]}. Its presence alone does not
establish a nonzero bias for a particular Gaussian mean model. Bias
reduction for nonlinear panels has an extensive literature {[}3, 4{]};
the calculation below concerns a fixed-grid Gaussian ODE mean with known
covariance, and uses a small-noise expansion rather than a large
within-record sample limit.

For one record whiten the observations by the known base covariance:
\(y=\mu(\theta_0,u)+\epsilon z\), \(z\sim N(0,I)\). At the true
parameters write \(S=\partial_\theta\mu\), \(Z=\partial_u\mu\),
\(N=Z^\top Z\) and \(P=I-ZN^{-1}Z^\top\). Assume the mean map is three
times continuously differentiable, \(N\) is nonsingular, and the
interior local nuisance minimizer admits a second-order expansion with
integrable, uniform remainders over the allowed initial states. Define

\[
h(u)=\sum_{i,j}(N^{-1})_{ij}\partial_{u_i u_j}\mu.
\]

\textbf{Local profile-score expansion.} For the unscaled least-squares
profile equation
\(\psi(\theta)=\partial_\theta\mu(\theta,\widehat u(\theta))^\top\{y-\mu(\theta,\widehat u(\theta))\}\),

\[
E\psi(\theta_0)=-\tfrac12\epsilon^2S^\top P h(u)+o(\epsilon^2).
\]

The log-likelihood score divides this expression by \(\epsilon^2\); the
order statement above concerns the unscaled equation. Suppose the
average local information
\(\overline I_e=\lim_n n^{-1}\sum_r S_r^\top P_r S_r\) is positive
definite, the population profile equation has a locally unique root, and
its uniform expansion is valid. That root has

\[
\theta_\epsilon^\dagger-\theta_0
=-\tfrac12\epsilon^2\overline I_e^{-1}
\lim_n n^{-1}\sum_r S_r^\top P_r h(u_r)+o(\epsilon^2).
\]

\textbf{Proof.} Write
\(\widehat u=u+\epsilon a+\epsilon^2b+o_{L^2}(\epsilon^2)\). The
nuisance normal equation gives \(a=N^{-1}Z^\top z\). The whitened
residual is
\(\epsilon Pz-\epsilon^2\{Zb+\tfrac12\mu_{uu}[a,a]\}+o_{L^1}(\epsilon^2)\).
Gaussian orthogonality makes \(a\) independent of \(Pz\), and
\(E(aa^\top)=N^{-1}\). Expanding the nuisance normal equation to second
order gives \(E b=-\tfrac12N^{-1}Z^\top h\). The mixed
coefficient-sensitivity term \(E\{(\mu_{\theta u}a)^\top Pz\}\) is zero.
Substitution into the coefficient equation therefore gives its stated
mean. Its population derivative at zero noise is \(-S^\top P S\), by
differentiating the nuisance normal equation. The local
implicit-function expansion after averaging records gives the displayed
root displacement. \(\square\)

A nonzero curvature coefficient predicts a fixed-noise displacement that
remains when the number of records increases. A consistency claim for
the computed local profile estimator additionally requires a uniform law
of large numbers and identification of this population root. Other
nuisance branches, boundary optima and rare large-noise fits fall
outside this local expansion. Linear mean models can have zero curvature
term. The experiment tests the size and direction of the displacement
rather than assuming a likelihood failure.

For the prescribed Van der Pol initial-state distribution, the
deterministic Kaggle calculation gives the coefficient of \(\epsilon^2\)
as approximately \((0.0024214,-0.0037071,0.0349238,-0.0152314)\). Eight-
and twelve-node Gauss--Legendre rules per initial-state coordinate agree
to relative \(1.8\times10^{-9}\); finite-difference checks of diagonal
second initial sensitivities have relative errors below
\(3.3\times10^{-6}\). At \(n=800\), the predicted relative bias divided
by the nuisance-adjusted norm standard error is 0.0803, 0.1605 and
0.3210 at 5\%, 10\% and 20\% noise. These are local predictions,
separate from the Monte Carlo pilot and the main comparison.

\subsection*{References}\label{references}
\addcontentsline{toc}{subsection}{References}

\protect\phantomsection\label{refs}
\begin{CSLReferences}{0}{0}
\bibitem[\citeproctext]{ref-neyman1948}
\CSLLeftMargin{{[}1{]} }%
\CSLRightInline{Neyman J, Scott EL. Consistent estimates based on
partially consistent observations. \emph{Econometrica} 16 (1948) 1--32.
\url{https://doi.org/10.2307/1914288}.}

\bibitem[\citeproctext]{ref-lancaster2000}
\CSLLeftMargin{{[}2{]} }%
\CSLRightInline{Lancaster T. The incidental parameter problem since
1948. \emph{Journal of Econometrics} 95 (2000) 391--413.
\url{https://doi.org/10.1016/S0304-4076(99)00044-5}.}

\bibitem[\citeproctext]{ref-hahn2004}
\CSLLeftMargin{{[}3{]} }%
\CSLRightInline{Hahn J, Newey WK. Jackknife and analytical bias
reduction for nonlinear panel models. \emph{Econometrica} 72 (2004)
1295--1319. \url{https://doi.org/10.1111/j.1468-0262.2004.00533.x}.}

\bibitem[\citeproctext]{ref-arellano2007}
\CSLLeftMargin{{[}4{]} }%
\CSLRightInline{Arellano M, Hahn J. Understanding bias in nonlinear
panel models: Some recent developments. \emph{Advances in economics and
econometrics: Theory and applications, ninth world congress} Cambridge
University Press 3 (2007) 381--409.
\url{https://doi.org/10.1017/CBO9780511607547.013}.}

\end{CSLReferences}

\end{document}
